\TieudeT{PHONG TỎA VẬT LÍ 12 - VẬT LÍ NHIỆT}
\subsubsection*{PHẦN I. Thí sinh trả lời câu hỏi từ câu 1 đến câu 18. Mỗi câu hỏi thí sinh chỉ chọn một phương án}
\setcounter{ex}{0}
\begin{ex}%Câu 1
	Phát biểu nào sau đây là đúng về nội năng của hệ khi hệ nhận nhiệt và nhận công?
	\choice
	{Chưa thể kết luận được}
	{Giảm}
	{Không đổi}
	{\True Tăng}
	\loigiai{
		\begin{itemize}
			\item $\Delta U = Q + A$ với $\begin{cases} Q>0\\ A>0 \end{cases} \Rightarrow \Delta U > 0$.
		\end{itemize}
	}
\end{ex}

\begin{ex}%Câu 2
		Phát biểu nào sau đây là đúng khi nói về nhiệt độ của chất lỏng? Trong quá trình sôi của chất lỏng 
	\choice
	{nhiệt độ luôn luôn tăng}
	{nhiệt độ luôn luôn giảm}
	{\True nhiệt độ luôn luôn không thay đổi}
	{nhiệt độ tăng hoặc giảm}
	\loigiai{
		
	}
\end{ex}

\begin{ex}%Câu 3
Trong quá trình sôi của chất lỏng \textbf{sai}?
	\choice
	{Nhiệt độ sôi của chất lỏng phụ thuộc vào áp suất trên bề mặt chất lỏng}
	{Trong quá trình sôi của chất lỏng thì nhiệt độ không thay đổi}
	{Mỗi chất lỏng sôi ở một nhiệt độ xác định và không đổi ở áp suất cho trước}
	{\True Sự bay hơi của chất lỏng chỉ xảy ra ở một nhiệt độ xác định}
	\loigiai{
		
	}
\end{ex}

\begin{ex}%Câu 4
	Phát biểu nào sau đây là đúng?
	\choice
	{Nội năng của một vật gồm tổng động năng của các phân tử cấu tạo nên vật}
	{Nội năng của một vật gồm tổng thế năng của các phân tử cấu tạo nên vật}
	{\True Nội năng của một vật gồm tổng động năng và thế năng của các phân tử cấu tạo nên vật}
	{Nội năng của một vật gồm tổng động năng và thế năng của vật}
	\loigiai{
		
	}
\end{ex}

\begin{ex}%Câu 5
	Khi khoảng cách giữa các phân tử rất nhỏ, thì giữa các phân tử
	\choice
	{chỉ có lực hút}
	{\True có cả lực hút và lực đẩy, nhưng lực đẩy lớn hơn lực hút}
	{chỉ có lực đẩy}
	{có cả lực hút và lực đẩy, nhưng lực đẩy nhỏ hơn lực hút}
	\loigiai{
		
	}
\end{ex}

\begin{ex}%Câu 6
	Phát biểu nào sau đây là đúng khi nói về điều kiện truyền nhiệt giữa hai vật?
	\choice
	{Nhiệt không thể truyền từ vật có nhiệt năng nhỏ sang vật có nhiệt năng lớn hơn}
	{Nhiệt không thể truyền giữa hai vật có nhiệt năng bằng nhau}
	{Nhiệt chỉ có thể truyền từ vật có nhiệt năng lớn hơn sang vật có nhiệt năng nhỏ hơn}
	{\True Nhiệt không thể tự truyền được từ vật có nhiệt độ thấp sang vật có nhiệt độ cao hơn}
	\loigiai{
	}
\end{ex}

\begin{ex}%Câu 7
	Nếu thả một cục nước đá có nhiệt độ $0^\circ\text{C}$ vào trong nước ở nhiệt độ $0^\circ\text{C}$ đựng trong một bình nhiệt lượng kế cách nhiệt tuyệt đối thì
	\choice
	{cục nước đá bị tan hoàn toàn thành nước}
	{cục nước đá bị tan một phần thành nước}
	{\True cục nước đá không tan trong nước}
	{cục nước đá làm một phần nước đông đặc thành nước đá}
	\loigiai{
		\begin{itemize}
			\item Nhiệt độ bằng nhau nên không có sự trao đổi nhiệt xảy ra.
		\end{itemize}
	}
\end{ex}

\begin{ex}%Câu 8
	Đặt cốc nhôm đựng $0,2$ lít nước ở nhiệt độ $30^\circ\text{C}$ (đo nhờ nhiệt kế $1-NK1$) vào trong bình cách nhiệt đựng $0,5$ lít nước ở nhiệt độ $60^\circ\text{C}$ (đo nhờ nhiệt kế $2-NK2$). Quan sát sự thay đổi nhiệt độ của nước trong bình và cốc từ khi bắt đầu thí nghiệm cho tới khi số chỉ hai nhiệt kế này bằng nhau. Nước trong bình truyền nhiệt năng cho nước trong cốc vì thấy số chỉ
	\choice
	{NK1 và số chỉ NK2 đều giảm}
	{NK1 và số chỉ NK2 đều tăng}
	{NK1 giảm còn số chỉ NK2 tăng}
	{\True NK1 tăng còn số chỉ NK2 giảm}
	\loigiai{
		\begin{itemize}
			\item Nước ở $30^\circ\text{C}$ tăng nhiệt độ, còn nước ở $60^\circ\text{C}$ giảm nhiệt độ. Chọn D
		\end{itemize}
	}
\end{ex}

\begin{ex}%Câu 9
	Hằng số $N_A=6,02.10^{23}\ mol^{-1}$. Số hạt helium có trong $0,5\ mol$ helium là
	\choice
	{$6,02.10^{23}$ hạt}
	{\True $3,01.10^{23}$ hạt}
	{$12,04.10^{23}$ hạt}
	{$24,08.10^{23}$ hạt}
	\loigiai{
		\begin{itemize}
			\item $N = n.N_A = 0,5.6,02.10^{23} = 3,01.10^{23}$ hạt.
		\end{itemize}
	}
\end{ex}

\begin{ex}%Câu 10
	Pha $100\ g$ nước ở $100^\circ\text{C}$ vào $100\ g$ nước ở $20^\circ\text{C}$ nhiệt độ cuối cùng của hỗn hợp nước là
	\choice
	{$100^\circ\text{C}$}
	{$20^\circ\text{C}$}
	{$50^\circ\text{C}$}
	{\True $60^\circ\text{C}$}
	\loigiai{
		\begin{itemize}
			\item Phương trình cân bằng nhiệt: $m_1c.(100-t) = m_2c.(t-20) \Rightarrow t = 60^\circ C$.
		\end{itemize}
	}
\end{ex}

\begin{ex}%Câu 11
	Pha $m$ (g) nước ở $100^\circ\text{C}$ vào $50\ g$ nước ở $30^\circ\text{C}$. Nhiệt độ cuối cùng của hỗn hợp nước là $50^\circ\text{C}$. Khối lượng $m$ là
	\choice
	{$10\ g$}
	{\True $20\ g$}
	{$30\ g$}
	{$40\ g$}
	\loigiai{
		\begin{itemize}
			\item $mc\Delta t = m'c\Delta t' \Rightarrow m.(100 - 50) = 50.(50 - 30) \Rightarrow m = 20\ g$.
		\end{itemize}
	}
\end{ex}

\begin{ex}%Câu 12
	Nhiệt độ của nước ở nhiệt độ phòng trong nhiệt giai Celsius là $28^\circ\text{C}$. Ứng với nhiệt giai Fahrenheit, nhiệt độ này là
	\choice
	{$15^\circ\text{F}$}
	{$48,6^\circ\text{F}$}
	{\True $82,4^\circ\text{F}$}
	{$47^\circ\text{F}$}
	\loigiai{
		\begin{itemize}
			\item $t_F = 1,8t_C + 32 = 1,8.28 + 32 = 82,4^\circ\text{F}$.
		\end{itemize}
	}
\end{ex}

\begin{ex}%Câu 13
	Đổ $200\ g$ nước ở nhiệt độ $t_1=20^\circ\text{C}$ vào $300\ g$ nước ở nhiệt độ $t_2$. Khi cân bằng nhiệt, nhiệt độ là $65^\circ\text{C}$. Nhiệt độ $t_2$ của nước là
	\choice
	{$42,5^\circ\text{C}$}
	{\True $95^\circ\text{C}$}
	{$85^\circ\text{C}$}
	{$45^\circ\text{C}$}
	\loigiai{
		\begin{itemize}
			\item $m_1c(t - t_1) = m_2c(t_2 - t) \Rightarrow 0,2.(65 - 20) = 0,3.(t_2 - 65) \Rightarrow t_2 = 95^\circ\text{C}$.
		\end{itemize}
	}
\end{ex}

\begin{ex}%Câu 14
	Một miếng nhôm nặng $1\ kg$ (nhôm có nhiệt dung riêng là $900\ J/kg.K$) được làm nóng để nhiệt độ của nó tăng thêm $5^\circ\text{C}$. Miếng nhôm đã nhận một nhiệt lượng bằng
	\choice
	{\True $4,5.10^3\ J$}
	{$9,0.10^3\ J$}
	{$1,4.10^4\ J$}
	{$2,0.10^4\ J$}
	\loigiai{
		\begin{itemize}
			\item $Q = mc\Delta t = 1.900.5 = 4500\ J$.
		\end{itemize}
	}
\end{ex}

\begin{ex}%Câu 15
	Biết nhiệt nóng chảy riêng của nước đá là $34.10^4\ J/kg$ và nhiệt dung riêng của nước là $4180\ J/kg.K$. Nhiệt lượng cần cung cấp để $4\ kg$ nước đá ở $0^\circ\text{C}$ để chuyển nó thành nước ở $20^\circ\text{C}$ là
	\choice
	{$1496400\ J$}
	{\True $1694400\ J$}
	{$1494600\ J$}
	{$1964400\ J$}
	\loigiai{
		\begin{itemize}
			\item $Q = m\lambda + mct = 4.34.10^4 + 4.4180.20 = 1694400\ J$.
		\end{itemize}
	}
\end{ex}

\begin{ex}%Câu 16
	Một viên đạn bằng bạc đang bay với vận tốc $200\ m/s$ thì va chạm vào một bức tường gỗ và nằm yên trong bức tường. Nhiệt dung riêng của bạc là $234\ J/(kg.K)$. Nếu coi viên đạn không trao đổi nhiệt với bên ngoài thì nhiệt độ của viên đạn sẽ tăng thêm 
	\choice
	{$58^\circ C$}
	{$171^\circ C$}
	{\True $85^\circ C$}
	{$250^\circ C$}
	\loigiai{
		\begin{itemize}
			\item $W = Q \Rightarrow \dfrac{mv^2}{2} = mc\Delta t \Rightarrow \Delta t = \dfrac{v^2}{2c} = \dfrac{200^2}{2.234} = 85,47^\circ C $.
		\end{itemize}
	}
\end{ex}

\begin{ex}%Câu 17
	Để xác định nhiệt nóng chảy của thiếc, người ta đổ $350\ g$ thiếc nóng chảy ở nhiệt độ $232^\circ\text{C}$ vào $330\ g$ nước ở $7^\circ\text{C}$ đựng trong một nhiệt lượng kế có nhiệt dung bằng $100\ J/K$. Sau khi cân bằng nhiệt, nhiệt độ của nước trong nhiệt lượng kế là $32^\circ\text{C}$. Biết nhiệt dung riêng của nước là $4,2\ J/g.K$, của thiếc rắn là $0,23\ J/g.K$. Nhiệt nóng chảy riêng của thiếc là
	\choice
	{$80\ J/g$}
	{\True $60\ J/g$}
	{$40\ J/g$}
	{$50\ J/g$}
	\loigiai{
		\begin{itemize}
			\item $m_t\lambda + m_tc_t(t_t - t) = (m_nc_n + C)(t - t_n)$  
			$\Rightarrow 350.\lambda + 350.0,23.(232 - 32) = (330.4,2 + 100).(32 - 7)$  
			$\lambda \approx 60\ J/g$.
		\end{itemize}
	}
\end{ex}
\begin{ex}%Câu 18
	Một nhiệt lượng kế bằng đồng thau khối lượng $128\ g$ chứa $210\ g$ nước ở nhiệt độ $8,4\ ^\circ C$. Người ta thả một miếng kim loại khối lượng $192\ g$ đã nung nóng tới $100^\circ C$ vào nhiệt lượng kế. Biết nhiệt độ khi bắt đầu có sự cân bằng nhiệt là $21,5\ ^\circ C$. Cho nhiệt dung riêng của nước là $4,18.10^3\ J/kgK$; của đồng thau là $0,128.10^3\ J/kgK$. Nhiệt dung riêng của chất làm miếng kim loại là
	\choice
	{$793,2\ J/kgK$}
	{$792,5\ J/kgK$}
	{$817,4\ J/kgK$}
	{\True $777,2\ J/kgK$}
	\loigiai{
	\begin{itemize}
		\item $\left( m_{Cu}.c_{Cu} + m_n.c_n\right) \Delta t = m_{kl}.c_{kl}.\Delta t_{kl} \Rightarrow c_{kl} \approx 777,2\ J/kg.K$.
	\end{itemize}
	
	}
\end{ex}
\subsubsection*{PHẦN 2. Thí sinh trả lời từ câu 1 đến câu 4. Trong mỗi ý a), b), c), d) ở mỗi câu, thí sinh chọn đúng hoặc sai.}
\setcounter{ex}{0}
\begin{ex}
	Khi nói về lực tương tác phân tử.
	\choiceTFt
	{\True Lực hút phân tử có thể có độ lớn bằng lực đẩy phân tử}
	{\True Lực hút phân tử có thể mạnh hơn lực đẩy phân tử}
	{Lực hút phân tử không thể yếu hơn lực đẩy phân tử}
	{\True Lực tương tác phân tử không đáng kể khi các phân tử ở xa nhau}
	\loigiai{
		
		\begin{itemchoice}
			\itemch 
			\itemch
			\itemch 
			\itemch
		\end{itemchoice}
	}
\end{ex}

\begin{ex}
	Một ấm điện có ghi $220\ V$ – $880\ W$ được mắc vào hiệu điện thế $220\ V$. Dùng ấm điện trên để đun sôi $2\ kg$ nước ở $25^\circ C$ trong $20$ phút. Biết nhiệt dung riêng của nước là $4200\ J/kgK$.
	\choiceTFt
	{\True Cường độ dòng điện qua ấm là $4\ A$}
	{\True Điện trở của ấm là $55\ \Omega$}
	{Nhiệt lượng có ích để đun sôi nước là $1056\ kJ$}
	{\True Hiệu suất của ấm điện là $59,66\%$}
	\loigiai{
		\begin{itemchoice}
			\itemch $ \mathcal{P} = UI \Rightarrow 880 = 220.I \Rightarrow I = 4\ A$.
			\itemch $ \mathcal{P} = \dfrac{U_{\text{đm}}^2}{R} \Rightarrow R = \dfrac{U_{\text{đm}}^2}{\mathcal{P}} = \dfrac{220^2}{880} = 55\ \Omega$.
			\itemch $Q_{ci} = m.c.\Delta t = 2.4200.75 = 630000\ J = 630\ kJ$.
			\itemch $\mathcal{H} = \dfrac{Q_{ci}}{Q} = \dfrac{630000}{880.20.60} \approx 0,5966 = 59,66\%$.
		\end{itemchoice}
	}
\end{ex}

\begin{ex}
	Một học sinh làm thí nghiệm như sau:\\
	Cho $1\ kg$ nước ở $10^\circ C$ rồi đặt lên bếp điện để đun đến khi nước sôi ở $100^\circ C$. Theo dõi thời gian đun học sinh đó ghi chép được các số liệu sau đây:
	\begin{itemize}
	\item Để đun nước từ $10^\circ C$ đến $100^\circ C$ cần $18$ phút.
	\item Kể từ lúc nước sôi thì sau $22$ phút thì lượng nước còn lại trong ấm là $0,8\ kg$.
	\end{itemize}
	 Biết nhiệt dung riêng của nước là $4180\ J/kg.K$ và công suất của bếp điện là không đổi trong quá trình đun, bỏ qua nhiệt dung của ấm và mất mát nhiệt ra môi trường.
	\choiceTFt
	{\True Nhiệt lượng cần để làm $1\ kg$ nước trong thí nghiệm trên sôi là $376200\ J$}
	{\True Nhiệt lượng cần để làm bay hơi $1\ kg$ nước ở $100^\circ C$ là $2299000\ J$}
	{\True Sau khi đun $62$ phút kể từ lúc đun lượng nước còn lại trong ấm là $0,6\ kg$}
	{Để cô cạn hoàn toàn nước trong ấm cần tổng thời gian đun là $110$ phút}
	\loigiai{
	\begin{itemchoice}
		\itemch $Q_a = mc\Delta t = 1.4180.90 = 376200\ J \Rightarrow \mathcal{P} = \dfrac{376200}{18.60} = \dfrac{1045}{3}\ W$.
		\itemch Ta có: $Q_b = m_h.L \Rightarrow \mathcal{P}.t_b = (1-0,8).L \Rightarrow \dfrac{1045}{3}.22.60 = 0,2.L \Rightarrow  L = 2299000\  J/kg $.
		\itemch $\mathcal{P}.(t_c-t_b) = m'_h.L \Rightarrow \dfrac{1045}{3}(62-18)60 = m'_h.L \Rightarrow  m'_h = 0,4\ kg \Rightarrow \Delta m = 1- 0,4 = 0,6\ kg$.
		\itemch Thời gian cô cạn: $ t = t_s + t_h = t_s + \dfrac{Q}{\mathcal{P}} = 18 + \dfrac{1.2299000}{\dfrac{1045}{3}.60} = 128$ phút.

	\end{itemchoice}
}
\end{ex}

\begin{ex}
	Một ô tô đã chạy quãng đường $6,9\ km$ với tốc độ không đổi trên một con đường thẳng, mất $5$ phút $45\ s$ và tiêu tốn $1,5\ kg$ nhiên liệu. Biết rằng lực kéo của ô tô là $2000\ N$ và năng suất tỏa nhiệt của nhiên liệu là $4,6.10^7\ J/kg$. Giả sử nhiên liệu cháy hết.
	\choiceTFt
	{Tốc độ của ô tô là $20\ km/h$}
	{\True Nhiệt lượng tỏa ra khi đốt cháy hoàn toàn nhiên liệu là $6,9.10^7\ J$}
	{\True Công suất lực kéo của ô tô là $40\ kW$}
	{\True Hiệu suất của động cơ ô tô là $20\ \%$}
	\loigiai{
		\begin{itemchoice}
			\itemch $v= \dfrac{s}{t} = \dfrac{6900}{5.60+45} = 20\ m/s = 72\ km/h$.
			\itemch $Q = m.q = 1,5.4,6.10^7 = 69.10^6\ J = 6,9.10^7\ J$.
			\itemch $\mathcal{P}_k = F.v = 2000.20 = 40000\ W = 40\ kW$.
			\itemch $\mathcal{H} = \dfrac{\mathcal{P}_k}{\mathcal{P}} = \dfrac{40000}{\dfrac{6,9.10^7}{5.60+45}} = 0,2 = 20\%$
		\end{itemchoice}
	}
\end{ex}
\subsubsection*{PHẦN 3. Thí sinh trả lời từ câu 1 đến câu 6.}
\setcounter{ex}{0}
\begin{ex}% Câu 1
	Nơi được ghi nhận là lạnh nhất thế giới mà vẫn có người sinh sống là vùng Oymyakon của đất nước Nga. Nơi đây từng ghi nhận nhiệt độ thấp nhất khoảng $-71^\circ C$. Ngược lại nơi nóng nhất thế giới mà vẫn có người sinh sống là thung lũng Chết ở Mĩ, nơi từng ghi nhận nhiệt độ cao nhất khoảng $53^\circ C$. Hỏi sự chênh lệch nhiệt độ cao nhất và thấp nhất ở hai địa điểm trên là bao nhiêu $^\circ C$?
	\SA[oly]{$124$}
	\loigiai{
		\begin{itemize}
			\item $\Delta t = 53-(-71) = 124^\circ C$.
		\end{itemize}
	}
\end{ex}

\begin{ex}% Câu 2
	Một bếp điện có ghi $220\ V$ – $1000\ W$ được sử dụng với hiệu điện thế $220\ V$ để đun sôi $2,5\ kg$ nước từ nhiệt độ ban đầu là $20\ ^\circ C$ thì mất thời gian là 14 phút 35 giây. Biết nhiệt dung riêng của nước là $c = 4200\ J/kgK$. Hiệu suất của bếp bằng bao nhiêu $\%$?
	\SA[oly]{$96$}
	\loigiai{
		\begin{itemize}
			\item $H = \dfrac{Q}{A} = \dfrac{mc\Delta t}{Pt} = \dfrac{2,5.4200.(100-20)}{1000.(14.60+35)} = 0,96 = 96\%$
		\end{itemize}
	}
\end{ex}

\begin{ex}% Câu 3
	Để biến $2\ kg$ nước ở $20^\circ C$ thành nước đá ở $0^\circ C$, người ta bỏ vào lượng nước trên một khối nước đá ở nhiệt độ $-10^\circ C$. Cho biết nhiệt dung riêng của nước là $c_{n} = 4200\ J/kg.K$, nhiệt dung riêng của nước đá là $c_{\text{đ}} = 2000\ J/kg.K$, nhiệt nóng chảy riêng của nước đá là $\lambda = 3,4.10^5\ J/kg$. Khối lượng nước đá tối thiểu cần dùng bằng bao nhiêu $kg$ (làm tròn kết quả đến chữ số hàng phần mười)?
	\SA[oly]{$42,4$}
	\loigiai{
		\begin{itemize}
			\item $m_n.c_n.\Delta t_n + m_n.\lambda = m_{\text{đ}}.c_{\text{đ}}.\Delta t_{\text{đ}} \Rightarrow 2.4200.20+2.3,4.10^5 = m_{\text{đ}}.2000.10 = 42,4\ kg$.
		\end{itemize}
	}
\end{ex}

\begin{ex}% Câu 4
	Để đo nhiệt độ trên bề mặt các thiên thể, các nhà thiên văn học thường dùng định luật dịch chuyển Wien (do nhà vật lí học người Đức Wilhelm Wien tìm ra): $\lambda_{max} = \dfrac{b}{T}$ với hằng số $b = 2898\  \mu m.K$ trong đó $\lambda_{max}$ là bước sóng ứng với cường độ lớn nhất. Mặt trời của chúng ta phát ra phần lớn bức xạ ở bước sóng $550,0\ nm$. Nếu một ngôi sao T phát ra phần lớn bức xạ ở bước sóng $137,5\ nm$ ứng với cường độ lớn nhất thì nhiệt độ tuyệt đối trên bề mặt của ngôi sao đó cao gấp bao nhiêu lần so với nhiệt độ tuyệt đối trên bề mặt của Mặt trời?
	\SA[oly]{4}
	\loigiai{
		\begin{itemize}
			\item $\lambda _1T_1 = \lambda _2T_2 = b \Rightarrow \dfrac{T_2}{T_1} = \dfrac{\lambda _1}{\lambda _2} = \dfrac{550}{137,5} = 4$.
		\end{itemize}
	}
\end{ex}

\noindent
\begin{minipage}[h]{0.55\textwidth}
	\textit{Sử dụng thông tin sau để trả lời câu 5 và câu 6:}
	
	Một học sinh tiến hành đun một khối nước đá ở điều kiện tiêu chuẩn từ $0^\circ C$ đến khi tan chảy hết thành nước và bay hơi ở $100^\circ C$. Hình vẽ là đồ thị biểu diễn sự phụ thuộc của nhiệt lượng mà khối nước đá nhận được từ lúc đun đến lúc bay hơi và sự thay đổi nhiệt độ của nó. Lấy nhiệt nóng chảy riêng của nước đá là $3,3. 10^5\ J/kg$, nhiệt dung riêng của nước gần đúng là $4200\ J/kg.K$, nhiệt hóa hơi của nước là $2,3.10^6\ J/kg$.
\end{minipage}
\begin{minipage}[h]{0.4\textwidth}
	\begin{tikzpicture}[draw=Mapcolor, very thick, >=latex']
		\draw [->](0,0) -- (6.5,0) node [above] {$Q\ (kJ)$};
		\draw[->] (0,0) -- (0,3) node [right] {$t\ ({^\circ}C)$};
		\draw [line width=1.5pt] (0,0) node [left] {$O$} 
		-- (1.65,0) circle (1pt) node [above] {$A$} 
		-- (3.75,2.5) circle (1pt) node [above] {$B$}  
		-- (6.05,2.5) circle (1pt) node [above] {$C$};
		\draw [dashed] (0,2.5) -- (3.75,2.5) -- (3.75,0) node [below] {$375$}
		node [below] at (1.65,0) {$165$}
		node [below] at (6.05,0) {$605$}
		node [left] at (0,2.5) {$100$};
	\end{tikzpicture}
\end{minipage}

\begin{ex}% Câu 5
	Khối lượng nước đá bằng bao nhiêu $g$?
	\SA[oly]{$500$}
	\loigiai{
		\begin{itemize}
			\item $ Q_A - 0 = m.\lambda \Rightarrow 165.10^3 = m.3,3.10^5 \Rightarrow m = 0,5\ kg = 500\ g$.
		\end{itemize}
	}
\end{ex}

\begin{ex}% Câu 6
	Khối lượng nước đã bay hơi trong giai đoạn $BC$ bằng bao nhiêu $g$?
	\SA[oly]{$100$}
	\loigiai{
		\begin{itemize}
			\item $Q_C-Q_B = m.L \Rightarrow (605 - 375).10^3 = m.2,3.10^6 \Rightarrow m = 0,1\ kg = 100\ g$.
		\end{itemize}
	}
\end{ex}
